% =====================================================================
\lezione{16/10/2026}{Laboratorio 2 -- Severity assessment e notifica del caso Sylvane Retraite}
  {Gianluca Gilardi}
  {svolgimento del Lab~2 (breach presso il fornitore del canale SFTP di un fondo pensione francese);
   worksheet di severity assessment; modulo di notifica alla CNIL; curveball}
\label{lez:2026-10-16}
% =====================================================================

\begin{nota}[Fonte]
Trascrizione automatica della registrazione: quasi tutta silenzio o rumore (all'inizio è stato
captato l'audio di un altro evento, non pertinente). Si riportano solo le indicazioni udibili di
Gianluca. Caso, metodologia e svolgimento ragionato nel Capitolo~\ref{lab:02}.
\end{nota}

\section{Sintesi}
Secondo laboratorio, condotto da Gianluca. I gruppi, nel ruolo di DPO del fondo pensione francese
Sylvane Retraite, compilano il \textbf{worksheet di severity assessment} (obbligatorio) e, se
concludono che la notifica è dovuta, il \textbf{modulo di notifica alla CNIL}. Durante il lavoro
arrivano nuove informazioni (\emph{curveball}) da integrare nella valutazione.

\section{Organizzazione}
\begin{itemize}
  \item Consegna sul repository SharePoint indicato nel caso: il solo worksheet, oppure un unico zip
    con worksheet e modulo di notifica se si notifica. Nomi del gruppo e dei componenti nel worksheet
    e nel modulo di caricamento; file nominati con il nome del gruppo.
  \item Gli ultimi cinque minuti sono stati lasciati per le domande; consegna entro le
    \textbf{13:30}.
  \item Lunedì~19/10: breve \textbf{debrief} sui lavori, poi si prosegue con la ISO.
\end{itemize}

\section{Indicazioni di Gianluca durante il lab}
\begin{esame}
Gianluca chiede una valutazione \textbf{complessiva} della violazione (severity e rischio
complessivi), non un punteggio per ogni singola voce dell'elenco dei dati: si valutano i singoli
elementi ``nella propria testa'' e si arriva a un \textbf{punteggio globale}, perché è questo che
chiede l'autorità di protezione dei dati.
\end{esame}
\begin{itemize}
  \item Le nuove informazioni sono state mostrate a schermo; l'ultima è stata annunciata come
    ``l'ultimo ostacolo'', dopo il quale i gruppi avevano tutti gli elementi disponibili.
\end{itemize}

\section{Domande aperte e punti da verificare}
\begin{verifica}
\begin{itemize}
  \item Nella cartella del lab è conservata solo la \textbf{curveball~1} (stessa comunicazione inviata
    a una società sorella che fa wealth management). Le altre informazioni mostrate a schermo, tra cui
    ``l'ultimo ostacolo'', non sono nel materiale né udibili: da recuperare (chat di Zoom o compagni).
  \item Contenuto del debrief del 19/10: da riportare nella voce di quella lezione.
\end{itemize}
\end{verifica}
